\documentclass[11pt,a4paper]{article}

% ---------------------------------------------------------------
%  LimpaC - Bases de Cálculo das Calculadoras de Impacto
% ---------------------------------------------------------------

\usepackage[T1]{fontenc}
\usepackage[utf8]{inputenc}
\usepackage[brazil]{babel}
\usepackage[a4paper,margin=2.2cm]{geometry}
\usepackage[table]{xcolor}
\usepackage{booktabs}
\usepackage{tabularx}
\usepackage{array}
\usepackage{amsmath}
\usepackage{siunitx}
\usepackage{enumitem}
\usepackage{tcolorbox}
\usepackage{titlesec}
\usepackage{fancyhdr}
\usepackage{graphicx}
\usepackage{hyperref}
\hypersetup{colorlinks=true,urlcolor={[HTML]{BE123C}},linkcolor={[HTML]{BE123C}}}
\tcbuselibrary{skins,breakable}

% ----------------------------- Paleta --------------------------
% Cores espelhadas do front-end (apps/web): brand rose + neutros slate.
%   --brand / --primary: rose   |  neutros: slate
\definecolor{limpacGreen}{HTML}{F43F5E}     % rose-500 (brand / accent)
\definecolor{limpacDark}{HTML}{0F172A}      % slate-900 (títulos)
\definecolor{limpacLight}{HTML}{FFF1F2}     % rose-50  (fundo de destaque)
\definecolor{limpacGray}{HTML}{64748B}      % slate-500 (texto auxiliar)
\definecolor{limpacBlue}{HTML}{334155}      % slate-700 (notas)
\definecolor{limpacBlueLight}{HTML}{F8FAFC} % slate-50  (fundo de notas)
\definecolor{roseDeep}{HTML}{BE123C}        % rose-700 (acento forte)

% ----------------------------- siunitx -------------------------
\sisetup{output-decimal-marker={,},group-separator={.},group-minimum-digits=4}

% ----------------------------- Títulos -------------------------
\titleformat{\section}
  {\color{limpacDark}\Large\bfseries\sffamily}
  {\thesection}{0.6em}{}
\titleformat{\subsection}
  {\color{limpacGreen}\large\bfseries\sffamily}
  {\thesubsection}{0.6em}{}
\titlespacing*{\section}{0pt}{1.4em}{0.6em}

% ----------------------------- Cabeçalho -----------------------
\pagestyle{fancy}
\fancyhf{}
\renewcommand{\headrulewidth}{0.4pt}
\renewcommand{\headrule}{\hbox to\headwidth{\color{limpacGreen}\leaders\hrule height \headrulewidth\hfill}}
\fancyhead[L]{\sffamily\small\color{limpacGray}LimpaC}
\fancyhead[R]{\sffamily\small\color{limpacGray}Bases de Cálculo das Calculadoras}
\fancyfoot[C]{\sffamily\small\color{limpacGray}\thepage}

\renewcommand{\arraystretch}{1.35}
\newcolumntype{Y}{>{\raggedright\arraybackslash}X}

% Caixa de fórmula reutilizável
\newtcolorbox{formulabox}[1][]{
  colback=limpacLight, colframe=limpacGreen,
  boxrule=0.8pt, arc=3pt, left=10pt, right=10pt, top=8pt, bottom=8pt,
  fonttitle=\bfseries\sffamily, #1}

\newtcolorbox{notebox}[1][]{
  colback=limpacBlueLight, colframe=limpacBlue,
  boxrule=0.6pt, arc=3pt, left=10pt, right=10pt, top=7pt, bottom=7pt, #1}

% ===============================================================
\begin{document}
\sffamily

% -------------------------- Capa -------------------------------
\begin{titlepage}
\centering
\vspace*{2.5cm}
{\color{limpacGreen}\rule{\linewidth}{2pt}}\\[0.6cm]
{\Huge\bfseries\color{limpacDark} LimpaC}\\[0.4cm]
{\LARGE\color{limpacGreen} Bases de Cálculo das Calculadoras de Impacto}\\[0.3cm]
{\large\color{limpacGray} Documentação técnica dos fatores ambientais e financeiros}\\[0.5cm]
{\color{limpacGreen}\rule{\linewidth}{2pt}}\\[1.2cm]

\begin{tcolorbox}[colback=limpacLight,colframe=limpacGreen,width=0.85\linewidth,
   arc=4pt,boxrule=1pt,halign=center]
\large
Este documento descreve, de forma prática, \textbf{como cada métrica é calculada},
quais \textbf{fatores unitários} são usados e quais \textbf{fontes} embasam esses números.
\end{tcolorbox}

\vfill
{\color{limpacGray}\large Documento técnico interno}\\[0.2cm]
{\color{limpacGray} Junho de 2026}
\end{titlepage}

% -------------------------- Sumário ----------------------------
\tableofcontents
\thispagestyle{fancy}
\vspace{1cm}

% ===============================================================
\section{Visão geral}

O LimpaC possui \textbf{duas calculadoras} independentes. Ambas seguem o mesmo
princípio: multiplicar uma \emph{quantidade informada pelo usuário} por um conjunto
de \emph{fatores unitários} (impacto por unidade evitada).

\begin{center}
\begin{tabularx}{\linewidth}{>{\bfseries\color{limpacDark}}p{4.2cm} Y}
\toprule
Calculadora & O que ela mede \\
\midrule
Cartões físicos & Impacto evitado ao substituir \textbf{cartões plásticos (PVC)} por
alternativas digitais. A entrada é o \textbf{número de cartões} não emitidos. \\
\addlinespace
Transações & Impacto evitado ao migrar \textbf{transações físicas} (dinheiro em espécie
+ recibo térmico) para \textbf{transações digitais}. A entrada é o \textbf{volume total}
de transações e o \textbf{percentual digital}. \\
\bottomrule
\end{tabularx}
\end{center}

\begin{notebox}
\textbf{Princípio de cálculo (modelo linear).} Todas as métricas são lineares:
\[
\text{Impacto} \;=\; \text{Quantidade evitada} \;\times\; \text{Fator unitário}.
\]
Os fatores são \emph{configuráveis} no backend
(\texttt{application.properties}), o que permite ajustar os valores sem alterar
o código-fonte.
\end{notebox}

% ===============================================================
\section{Calculadora de Cartões Físicos}

\subsection{Entrada e modelo}

A única entrada é a quantidade de cartões físicos evitados, $N_{\text{cartões}}$.
Cada métrica $M$ é obtida por:
\[
M \;=\; N_{\text{cartões}} \times f_M
\]
onde $f_M$ é o fator unitário da tabela abaixo. O número de árvores é
\textbf{arredondado para o inteiro mais próximo}; as demais métricas mantêm casas decimais.

\subsection{Fatores unitários (por cartão evitado)}

\begin{center}
\begin{tabularx}{\linewidth}{Y c c}
\toprule
\rowcolor{limpacLight}
\textbf{Métrica} & \textbf{Fator por cartão} & \textbf{Unidade} \\
\midrule
CO\textsubscript{2} evitado          & \num{0.0471}  & kg \\
Plástico (PVC) evitado               & \num{0.00714} & kg \\
Árvores preservadas                  & \num{0.0342}  & árvore \\
Água economizada                     & \num{12.857}  & litros \\
Energia economizada                  & \num{0.514}   & kWh \\
Dinheiro economizado                 & \num{2.77}    & R\$ \\
\bottomrule
\end{tabularx}
\end{center}

\subsection{Composição do valor financeiro}

O valor economizado por cartão (\textbf{R\$ 2,77}) é a soma de três parcelas,
calculadas separadamente no backend para facilitar auditoria:

\begin{center}
\begin{tabularx}{\linewidth}{Y c}
\toprule
\rowcolor{limpacLight}
\textbf{Parcela} & \textbf{Valor (R\$/cartão)} \\
\midrule
Material (blank PVC CR80)                  & \num{0.49} \\
Fabricação / personalização / impressão    & \num{1.58} \\
Logística / frete (rateado por lote)       & \num{0.70} \\
\midrule
\textbf{Total}                             & \textbf{\num{2.77}} \\
\bottomrule
\end{tabularx}
\end{center}

\begin{formulabox}[title=Exemplo prático --- 350 cartões evitados]
\begin{tabular}{@{}ll@{}}
CO\textsubscript{2}:    & $350 \times 0{,}0471 = \mathbf{16{,}5}$ kg \\
Plástico:               & $350 \times 0{,}00714 = \mathbf{2{,}50}$ kg \\
Árvores:                & $350 \times 0{,}0342 = 11{,}97 \approx \mathbf{12}$ árvores \\
Água:                   & $350 \times 12{,}857 = \mathbf{4.500}$ L \\
Energia:                & $350 \times 0{,}514 = \mathbf{179{,}9}$ kWh \\
Economia:               & $350 \times 2{,}77 = \mathbf{R\$\ 969{,}50}$ \\
\end{tabular}
\end{formulabox}

\subsection{Fontes e premissas}
\begin{itemize}[leftmargin=1.4em]
\item \textbf{Peso / plástico:} um cartão de pagamento PVC pesa em média
$\approx \SI{5}{g}$; o fator usado ($\approx \SI{7}{g}$) inclui uma margem para
laminação e acabamento.
\item \textbf{CO\textsubscript{2}:} a produção de PVC tem GWP de
$\approx \SI{1.88}{kg\,CO_2/kg}$. O fator de \SI{47}{g/cartão} corresponde ao
material PVC mais um rateio de frete rodoviário, sendo um valor \emph{conservador}
frente às estimativas de ciclo de vida completo (150--170 g/cartão).
\item \textbf{Custo:} blank PVC CR80 em volume gira em torno de
R\$ 0,49--0,56/un.; cartões personalizados de baixo volume amortizam
fabricação em torno de R\$ 1,50--2,00/un.
\end{itemize}

% ===============================================================
\section{Calculadora de Transações}

\subsection{Entrada e modelo}

A calculadora recebe o \textbf{volume total de transações} ($T$) e o
\textbf{percentual digital} ($p$, em \%). A partir disso deriva:
\[
T_{\text{digital}} = T \times \frac{p}{100}
\qquad\qquad
T_{\text{física}} = T - T_{\text{digital}}
\]

\begin{notebox}
\textbf{Importante:} o impacto evitado é calculado \textbf{somente sobre as
transações digitais} ($T_{\text{digital}}$), pois representam aquelas que
deixaram de usar dinheiro em espécie + recibo térmico.
\end{notebox}

Cada métrica $M$ segue:
\[
M \;=\; T_{\text{digital}} \times f_M
\]

\subsection{Fatores unitários (por transação digital)}

\begin{center}
\begin{tabularx}{\linewidth}{Y c c}
\toprule
\rowcolor{limpacLight}
\textbf{Métrica} & \textbf{Fator por transação} & \textbf{Unidade} \\
\midrule
CO\textsubscript{2} evitado    & \num{0.007}     & kg \\
Papel (recibo térmico) evitado & \num{0.0023}    & kg \\
Água economizada               & \num{0.133}     & litros \\
Árvores preservadas            & \num{0.0000117} & árvore \\
Dinheiro economizado           & \num{0.34}      & R\$ \\
\bottomrule
\end{tabularx}
\end{center}

\subsection{Composição do valor financeiro}

\begin{center}
\begin{tabularx}{\linewidth}{Y c}
\toprule
\rowcolor{limpacLight}
\textbf{Parcela} & \textbf{Valor (R\$/transação)} \\
\midrule
Recibo térmico (rolo de papel rateado)        & \num{0.06} \\
Manuseio de numerário (contagem, conferência, & \\
\quad transporte de valores --- rateado)      & \num{0.28} \\
\midrule
\textbf{Total}                                & \textbf{\num{0.34}} \\
\bottomrule
\end{tabularx}
\end{center}

\begin{formulabox}[title={Exemplo prático --- 10.000 transações, 70 por cento digitais}]
$T_{\text{digital}} = 10.000 \times 0{,}70 = 7.000$ transações digitais.
\begin{tabular}{@{}ll@{}}
CO\textsubscript{2}:  & $7.000 \times 0{,}007 = \mathbf{49}$ kg \\
Papel:                & $7.000 \times 0{,}0023 = \mathbf{16{,}1}$ kg \\
Água:                 & $7.000 \times 0{,}133 = \mathbf{931}$ L \\
Árvores:              & $7.000 \times 0{,}0000117 = \mathbf{0{,}08}$ árvore \\
Economia:             & $7.000 \times 0{,}34 = \mathbf{R\$\ 2.380{,}00}$ \\
\end{tabular}
\end{formulabox}

\subsection{Fontes e premissas}
\begin{itemize}[leftmargin=1.4em]
\item \textbf{CO\textsubscript{2}:} relatório \emph{Skip the Slip}
(Green America / NRDC): $\approx$ 4 bilhões de libras de CO\textsubscript{2}/ano
sobre $\approx$ 256 bilhões de recibos nos EUA $\Rightarrow \approx \SI{7}{g/recibo}$.
\item \textbf{Papel:} um recibo térmico médio pesa 2--3 g $\Rightarrow \SI{0.0023}{kg}$.
\item \textbf{Água:} $\approx$ 9 bilhões de galões/ano $\div$ 256 bilhões de recibos
$\Rightarrow \approx \SI{0.13}{L/recibo}$.
\item \textbf{Árvores:} $\approx$ 3 milhões de árvores/ano $\div$ 256 bilhões de
recibos $\Rightarrow \approx \num{0.0000117}$ árvore/recibo.
\item \textbf{Manuseio de numerário:} proxy conservador informado por dados de
perdas no varejo (Abrappe 2024, $\approx$ 1,51\% da receita líquida), amortizado
por transação.
\end{itemize}

% ===============================================================
\section{Resumo dos fatores}

\begin{center}
\begin{tabularx}{\linewidth}{Y c c}
\toprule
\rowcolor{limpacDark}
\textbf{\color{white}Fator} & \textbf{\color{white}Cartão} & \textbf{\color{white}Transação} \\
\midrule
CO\textsubscript{2} (kg)        & \num{0.0471}  & \num{0.007} \\
\rowcolor{limpacLight}
Plástico / Papel (kg)           & \num{0.00714} & \num{0.0023} \\
Água (L)                        & \num{12.857}  & \num{0.133} \\
\rowcolor{limpacLight}
Árvores                         & \num{0.0342}  & \num{0.0000117} \\
Energia (kWh)                   & \num{0.514}   & --- \\
\rowcolor{limpacLight}
Economia (R\$)                  & \num{2.77}    & \num{0.34} \\
\bottomrule
\end{tabularx}
\end{center}

\vspace{0.4cm}
\begin{notebox}
\textbf{Onde alterar.} Todos os fatores vivem em
\texttt{apps/backend/src/main/resources/application.properties}, sob os prefixos
\texttt{app.metrics.*} (cartões) e \texttt{app.transaction-metrics.*} (transações).
Os valores financeiros são somados em tempo de execução a partir de suas parcelas.
\end{notebox}

% ===============================================================
\section{Fontes}

Referências utilizadas para embasar os fatores unitários. Os valores adotados no
LimpaC são \emph{proxies conservadores} derivados destas fontes.

\subsection{Cartões físicos (PVC)}
\begin{enumerate}[leftmargin=1.6em]
\item Green America --- \emph{Skip the Slip} (programa de referência sobre
impacto de materiais descartáveis):\\
\href{https://greenamerica.org/save-trees/skip-slip}{greenamerica.org/save-trees/skip-slip}
\item Design Life-Cycle --- \emph{Credit/Debit Card} (peso $\approx$ 5 g,
composição PVC e ciclo de vida):\\
\href{http://www.designlife-cycle.com/creditdebit-card}{designlife-cycle.com/creditdebit-card}
\item CleverCards --- \emph{The climate crisis behind plastic cards}
(150--170 g CO\textsubscript{2}/cartão no ciclo completo):\\
\href{https://www.clevercards.com/blog/the-climate-crisis-behind-plastic-cards}{clevercards.com/blog/the-climate-crisis-behind-plastic-cards}
\item ICMA --- \emph{Eco-Friendly Materials in Card Manufacturing}
(materiais, custos e alternativas recicladas):\\
\href{https://icma.com/eco-friendly-materials-in-card-manufacturing}{icma.com/eco-friendly-materials-in-card-manufacturing}
\item GWP do PVC ($\approx \SI{1.88}{kg\,CO_2/kg}$) e custos de blank CR80:
referências de mercado de impressão de cartões (Perfect Plastic Printing,
\href{https://perfectplastic.com/carbon-neutral-footprint/}{perfectplastic.com/carbon-neutral-footprint}).
\end{enumerate}

\subsection{Transações (recibo térmico + numerário)}
\begin{enumerate}[leftmargin=1.6em]
\item Green America / NRDC --- \emph{Skip the Slip} report (4 bi lb
CO\textsubscript{2}, 3 mi árvores, 9 bi galões de água sobre $\approx$ 256 bi
recibos/ano nos EUA):\\
\href{https://greenamerica.org/press-release/skip-slip-report-toxic-paper-receipts-jeopardize-health}{greenamerica.org/press-release/skip-slip-report-toxic-paper-receipts}
\item Climate Action --- \emph{Report: Skip the Slip} (custos ambientais e
riscos à saúde dos recibos de papel):\\
\href{https://www.climateaction.org/news/report-skip-the-slip-environmental-costs-human-health-risks-of-paper-receip}{climateaction.org/news/report-skip-the-slip}
\item Ingenico --- \emph{Why digital receipts could be the perfect first step}
(estimativa de CO\textsubscript{2} por recibo):\\
\href{https://ingenico.com/en/newsroom/blogs/sustainability-payments-digital-receipts}{ingenico.com/.../sustainability-payments-digital-receipts}
\item eyos retail --- \emph{The Environmental Cost of Paper Receipts}
(peso médio do recibo térmico, 2--3 g):\\
\href{https://eyos.one/blog/the-environmental-cost-of-paper-receipts/}{eyos.one/blog/the-environmental-cost-of-paper-receipts}
\item Abrappe --- dados de perdas no varejo 2024 ($\approx$ 1,51\% da receita
líquida), usados como proxy do custo de manuseio de numerário.
\end{enumerate}

\begin{notebox}
\textbf{Nota de calibração.} Onde uma fonte traz faixa ou valor de ciclo de vida
completo (ex.: 150--170 g CO\textsubscript{2}/cartão), o LimpaC adota a parcela
\emph{material + logística} (\SI{47}{g}), mais conservadora e auditável. Os
fatores podem ser recalibrados a qualquer momento em \texttt{application.properties}.
\end{notebox}

\end{document}
